\newpage
Let $\alpha: R \to S$ be a ring homomorphism
such that $S$ is a free-$R$-module of rank $n<\infty$.
Let $M, P$ be $S$-modules.
Then we have the corestriction map
\(
    \Cores_R^S: \ext_R^i(P, M) \to \ext_S^i(P, M)
\)
and the restriction map
\(
    \Res_R^S: \ext_S^i(P, M) \to \ext_R^i(P, M).
\)

\begin{proposition}
    The composition $\Cores_R^S \circ \Res_R^S$ is multiplication by $n$.
\end{proposition}
\begin{proof}
    First, we start with the case $i=0$.
    We have the following diagram:
    \[
        \begin{tikzcd}
            \Hom_S(P, M) \arrow[r, hook, "\Res_R^S"] & \Hom_R(P, M) \arrow[r, "\sim"] &\Hom_S(P\underset{R}{\otimes}S, M)\arrow[d, "\sim" sloped]\\
            \Hom_S(P, M) &\arrow[l, "\Delta^*"'] \bigoplus \Hom_S(Pe_i, M) & \arrow[l, "\sim"'] \Hom_S(P^n, M)
        \end{tikzcd}
    \]
    where $e_1, \ldots, e_n$ is a basis of $S$ as a free-$R$-module.
    We have $\Delta^*(f_1, \ldots, f_n) = \sum_{i=1}^n f_i$.
    We observe that the composition gives $f \mapsto (f, \ldots, f) \mapsto \sum_{i=1}^n f = nf$.

    It gives a natural transformation $\eta: \Hom_S(P, -) \to \Hom_S(P, -)$
    on $S$-$\Mod$ such that $\eta(M)$ is multiplication by $n$.
    Since $\Hom_S(P, -)$ is a left exact functor, we can extend $\eta$ to a natural transformation
    \(
        \eta^i: \ext_S^i(P, -) \to \ext_S^i(P, -)
    \)
    for all $i \ge 0$, and it must be multiplication by $n$ as well.
\end{proof}

\begin{corollary}
    Let $G$ be a discrete group and $H$ be a subgroup of a finite
    index $n$.
    Let $A \in G\text{-}\Mod$.
    Then the composition
    \begin{align*}
        H^i(G, A) \xrightarrow{\Res_H^G} H^i(H, A) \xrightarrow{\Cores_H^G} H^i(G, A)
    \end{align*}
    is multiplication by $n = |G/H|$.
\end{corollary}

This can be interpreted as the algebraic version of a topological fact.
Let $X$ be a connected topological space and $Y\to X$ be an $n$-sheeted covering space.
Then the restriction map is corresponding to the pullback $p^*$ and
the corestriction map is corresponding to the transfer map $p_!$.
(So we sometimes call the corestriction map the transfer map.)
Copying a class to $n$ sheets and adding them back up gives $n$ times the original class.

\begin{corollary}
    Suppose $G$ is a finite group of order $n$
    and $A$ is a $G$-module.
    Then $H^i(G, A)$ is annihilated by $n$ for all $i>0$.
\end{corollary}
\begin{proof}
    Apply the previous corollary with the trivial subgroup $H = \{e\}$.
    Then we observe $H^i(G, A) \to H^i(\{e\}, A) \to H^i(G, A)$ is multiplication by $n$.
    But $H^i(\{e\}, A) = 0$ for all $i>0$, so we conclude that $H^i(G, A)$ is annihilated by $n$ for all $i>0$.
\end{proof}

Let $G$ be a discrete group and $H$ be a normal subgroup of finite index,
and $A$ be a $G$-module.
Then we have the inflation map, the